% Netzwerk-Spiel: gemeinsame Daten und Makros für die Druckvorlagen
% (Netzwerkspiel_Gruppen.tex und Netzwerkspiel_Pakete.tex)

% Schichtfarben wie in den Slides
\providecommand{\applicationLayerColor}{Gold}
\providecommand{\transportLayerColor}{Indigo}
\providecommand{\networkLayerColor}{RoyalBlue}
\providecommand{\dataLinkLayerColor}{Gray}

% Anzahl Gruppen (je 1 Router und 3 oder 4 Hosts), höchstens 8.
% n SuS: floor(n/4) Gruppen, davon (n mod 4) Gruppen mit 5 Personen.
\providecommand{\nsAnzahlGruppen}{6}

% Verkürzte MAC-Adressen (3 statt 6 Bytes), einmal zufällig erzeugt und fest.
% Gerät 0 = Router, 1 bis 4 = Hosts (Host 4 nur in Fünfergruppen).
\csdef{nsmac-1-0}{A5:4D:CA} \csdef{nsmac-1-1}{18:25:30} \csdef{nsmac-1-2}{BB:1D:6D} \csdef{nsmac-1-3}{13:2C:DE} \csdef{nsmac-1-4}{D6:23:7B}
\csdef{nsmac-2-0}{2E:D9:1E} \csdef{nsmac-2-1}{3F:72:1F} \csdef{nsmac-2-2}{CB:19:71} \csdef{nsmac-2-3}{17:44:94} \csdef{nsmac-2-4}{D6:49:3C}
\csdef{nsmac-3-0}{9D:5C:34} \csdef{nsmac-3-1}{60:BE:31} \csdef{nsmac-3-2}{20:1E:69} \csdef{nsmac-3-3}{FE:DA:A0} \csdef{nsmac-3-4}{EE:E8:B9}
\csdef{nsmac-4-0}{99:7F:5C} \csdef{nsmac-4-1}{7C:29:99} \csdef{nsmac-4-2}{FD:AF:E5} \csdef{nsmac-4-3}{93:25:3C} \csdef{nsmac-4-4}{D6:54:AF}
\csdef{nsmac-5-0}{4D:FA:D7} \csdef{nsmac-5-1}{14:27:A0} \csdef{nsmac-5-2}{AE:B3:FE} \csdef{nsmac-5-3}{E9:23:2F} \csdef{nsmac-5-4}{8A:F2:21}
\csdef{nsmac-6-0}{1F:9E:E4} \csdef{nsmac-6-1}{91:C5:B1} \csdef{nsmac-6-2}{0B:EC:B5} \csdef{nsmac-6-3}{56:3B:FC} \csdef{nsmac-6-4}{1E:6F:93}
\csdef{nsmac-7-0}{42:7E:CB} \csdef{nsmac-7-1}{C8:FE:29} \csdef{nsmac-7-2}{55:E5:CD} \csdef{nsmac-7-3}{8E:46:DC} \csdef{nsmac-7-4}{8E:D4:B7}
\csdef{nsmac-8-0}{C2:76:4D} \csdef{nsmac-8-1}{2A:5A:4D} \csdef{nsmac-8-2}{76:77:06} \csdef{nsmac-8-3}{F8:5D:86} \csdef{nsmac-8-4}{90:02:4A}

\newcommand{\nsMac}[2]{\texttt{\csuse{nsmac-#1-#2}}}

% Partnergruppe für Runde 3: 1-2, 3-4, ...; bei ungerader Anzahl spielt die
% letzte Gruppe mit der vorletzten, die dann zwei Nachbarn hat (\nsExtra).
\newcommand{\nsPartner}[1]{\the\numexpr #1\ifodd#1 \ifnum#1=\nsAnzahlGruppen\relax -1\else +1\fi\else -1\fi\relax}
\newcommand{\nsHatExtra}[1]{\ifodd\nsAnzahlGruppen \ifnum\numexpr#1+1\relax=\nsAnzahlGruppen\relax \ifnum\nsAnzahlGruppen>1 1\else 0\fi\else 0\fi\else 0\fi}
\newcommand{\nsExtra}{\nsAnzahlGruppen}

% Ganze Seite als TikZ-Overlay; Koordinaten in mm ab linker unterer Ecke
\newcommand{\nsSeite}[1]{%
	\clearpage\thispagestyle{empty}\null
	\begin{tikzpicture}[remember picture,overlay,x=1mm,y=1mm,shift={(current page.south west)}]
		#1
	\end{tikzpicture}%
	\clearpage
}
\tikzset{
	nsschnitt/.style={dashed,gray,thin},
	nsfalz/.style={dotted,gray,thick},
	nsfeld/.style={draw=black!60,thin},
}

% Kopfband: #1 Farbe, #2 x0, #3 y0, #4 x1, #5 y1, #6 Text
\newcommand{\nsBand}[6]{%
	\fill[#1!20] (#2,#3) rectangle (#4,#5);
	\node[anchor=west,font=\bfseries] at (#2+2,{(#3+#5)/2}) {#6};
}

% Ethernet-Frame (A4, einmal falten): #1 Gruppe, #2 Ziel-MAC, #3 Bemerkung
\newcommand{\nsFrame}[3]{%
	\nsSeite{
		\nsBand{\dataLinkLayerColor}{12}{265}{198}{285}{\Large Ethernet-Frame (Netzzugangsschicht)}
		\node[anchor=east,font=\large] at (196,275) {Gruppe #1};
		\node[font=\Large] at (105,243) {Ziel-MAC};
		\node[font=\fontsize{58}{60}\selectfont\ttfamily\bfseries] at (105,218) {#2};
		\node[font=\Large] at (105,188) {Absender-MAC (Router Gruppe #1): \nsMac{#1}{0}};
		\node[font=\small,gray] at (105,172) {#3};
		\draw[nsfalz] (0,148.5) -- (210,148.5);
		\node[fill=white,font=\footnotesize,text=gray] at (105,148.5) {hier falten, Aufdruck aussen};
		\node[font=\Large,align=center] at (105,115) {Hier das IP-Paket hineinlegen.};
		\node[font=\normalsize,align=center,text width=160mm] at (105,40) {Empfänger: IP-Paket herausnehmen und das leere Frame-Blatt dem Router zurückgeben. Der Router verwendet es wieder.};
	}
}

% IP-Kärtchen (12 pro A4-Seite, 105 x 49.5 mm): #1 x0, #2 y0
\newcommand{\nsIpPaket}[2]{%
	\nsBand{\networkLayerColor}{#1+5}{#2+40}{#1+100}{#2+47}{\small IP-Paket (Internetschicht)}
	\node[anchor=west] at (#1+7,#2+32) {Von-IP:};
	\node[anchor=west,font=\ttfamily] at (#1+27,#2+32) {192.168.\rule{9mm}{.3pt}.\rule{9mm}{.3pt}};
	\node[anchor=west] at (#1+7,#2+21) {An-IP:};
	\node[anchor=west,font=\ttfamily] at (#1+27,#2+21) {192.168.\rule{9mm}{.3pt}.\rule{9mm}{.3pt}};
	\node[anchor=west] at (#1+7,#2+9) {TTL:};
	\draw[nsfeld] (#1+27,#2+5) rectangle (#1+42,#2+13);
	\node[anchor=west,font=\scriptsize,gray] at (#1+44,#2+9) {nur Runde 3, Start: 4};
}

% Nachrichtenzettel (16 pro A4-Seite, 105 x 37.1 mm): #1 x0, #2 y0
\newcommand{\nsZettel}[2]{%
	\fill[\transportLayerColor!20] (#1+5,#2+27) rectangle (#1+100,#2+34);
	\node[anchor=west,font=\small] at (#1+7,#2+30.5) {\textbf{Seq} \rule{8mm}{.3pt} von \rule{8mm}{.3pt}};
	\node[anchor=east,font=\scriptsize] at (#1+98,#2+30.5) {Transportschicht};
	\fill[\applicationLayerColor!20] (#1+5,#2+5) rectangle (#1+100,#2+25);
	\foreach \i in {0,...,9} {\draw[nsfeld,fill=white] (#1+10+8.5*\i,#2+8) rectangle (#1+18.5+8.5*\i,#2+20);}
	\node[font=\tiny] at (#1+52.5,#2+22.5) {Nachricht (Anwendungsschicht): ein Zeichen pro Feld, Leerzeichen zählen mit};
}

% Gerätekarte (A7): #1 x0, #2 y0, #3 Gruppe, #4 Gerätenummer (0 = Router)
\newcommand{\nsKarte}[4]{%
	\ifnum#4=0
		\nsBand{\dataLinkLayerColor}{#1+5}{#2+60}{#1+100}{#2+70}{Gruppe #3, Router}
	\else
		\nsBand{\dataLinkLayerColor}{#1+5}{#2+60}{#1+100}{#2+70}{Gruppe #3, Computer\ifnum#4=4\ (nur Fünfergruppe)\fi}
	\fi
	\node[anchor=west] at (#1+7,#2+51) {Name: \rule{65mm}{.3pt}};
	\node[anchor=west] at (#1+7,#2+38) {MAC:};
	\node[anchor=west,font=\LARGE\ttfamily\bfseries] at (#1+22,#2+38) {\csuse{nsmac-#3-#4}};
	\ifnum#4=0
		\node[anchor=west] at (#1+7,#2+25) {IP: \texttt{192.168.#3.1}};
	\else
		\node[anchor=west] at (#1+7,#2+25) {IP: \rule{50mm}{.3pt}};
		\node[anchor=west,font=\small,gray] at (#1+7,#2+14) {IP-Adresse beim Router holen};
	\fi
}
